\section{Further research questions and proposed milestones}
\label{sec:research}

The next stage should pursue two complementary goals: improve the exact
implemented recognizer on difficult inputs, and prove structural statements
strong enough to control its worst case.  The following questions specify
what would count as progress.  They include possible negative results; a
counterexample to an attractive compression hypothesis is useful when it
prevents an unjustified global claim.

\subsection{Q1. Which frontier topology controls matching support?}

The immediate question is whether a useful intermediate bound lies between
the single-disk Catalan estimate and the unrestricted double factorial.
Suppose the processed region has $b$ relevant boundary components.  Can its
realizable matching support be bounded, under a precise encoding of endpoint
orders and boundary incidences, by an expression such as
$2^{O(w)}w^{O(b)}$?  If so, which additional winding or connectivity data
must be retained?  If not, can an explicit planar knot family demonstrate
the obstruction?

The first milestone is a certified boundary representation derived from
the PD rotation system.  The checker should report actual boundary cycles,
the frontier incidences on each cycle, and the processed side.  A second
milestone is a theorem about exactly that representation, tested against
the six-state counterexample and the default Conway-sum fixture.  Counting
abstract noncrossing partitions under a boundary order not supplied by the
input would not meet this milestone.

\subsection{Q2. Can certified noose decompositions replace greedy scans?}

A greedy linear order is easy to run and often effective, but it lacks the
geometric certificate required for the sharper count.  A decomposition based
on embedded separators or nooses could carry that certificate with each
interface.  The design question is how to compose tangle complexes along a
branching decomposition while controlling both width and homological
multiplicity.

The first prototype should expose a checked decomposition format before
optimizing the algebra.  Compare its certified boundary widths, support
sizes, peak multiplicities, and wall-clock cost with the existing greedy
order on identical diagrams.  A low-width decomposition can still lose
through expensive joining operations; the accounting must include those
joins.  A route to \eqref{eq:main-target} needs a structural theorem after
allowed reductions, or a complementary method for pieces that retain large
interfaces.

\subsection{Q3. What compresses repeated tangle complexes?}

Visible two-edge cuts admit a scalar connected-sum rule and are now handled
efficiently.  Four-ended cuts are the next natural extension.  Can a useful
class of tangles be represented by small complexes of boundary modules,
with exact composition and a canonical compressed form?  Repeated tangles
could then share a symbolic representation instead of duplicating all
generators.

The central obstacle is extension data.  The total rank, graded rank, or
Jones polynomial of each piece generally does not specify the differential
or the boundary action needed after gluing.  A proposed summary must therefore
come with a congruence theorem: equal summaries give equivalent results under
every permitted complementary gluing.  Test this property on exhaustive
small tangles with several closures before treating it as an algorithmic
interface.  A compact dictionary of repeated byte-identical complexes can
be an initial engineering step, but is not itself such a theorem.

\subsection{Q4. Can decision certificates appear before full reduction?}

The isolated-summand theorem supplies one exact lower bound, and the
vertex-cover lemma supplies another after scalar closure.  Can they be
extended to small blocks of objects whose every allowed completion has
nonzero reduced homology?  A useful block certificate must survive the
remaining differential and gluing functors, rather than merely count
uncancelled generators.

Implement the scalar support-graph cover bound first and measure its overhead
against the final cancellation it avoids.  Then study small open complexes
with certified nonzero closures.  A compelling positive result would be a
class of prime-looking filter survivors where a bounded-size certificate
arises after only a controlled portion of the scan.  A compelling negative
result would show that an entire proposed class of local certificates can
remain absent until the end despite large eventual rank.  Neither outcome
requires pretending that the present connected-sum examples establish new
pipeline superiority.

\subsection{Q5. Can the pointed backend be selected cheaply?}

The last-crossing edge policy has a proved frontier property, but it is not
an optimal basepoint theorem.  Which inexpensive statistics predict when
the quotient's savings exceed its bookkeeping?  Candidate statistics include
the time of first marked-edge appearance, early matching multiplicities,
the distribution of equal-matching pivots, and the number of marked surface
components in compiled plans.

A selection rule should be trained or tuned on one corpus and evaluated on
a separate corpus.  Its entire cost belongs in the end-to-end timing.
Running both expensive scans merely to choose the faster one is not a
low-cost selector, although a deliberately bounded race may be useful under
a stated resource policy.  Further algebra work could seek a fused ordinary
and pointed transfer routine that shares topology compilation without mixing
the incompatible numeric cache bases.

\subsection{Q6. Which Alexander matrices have controlled elimination work?}

The sparse theorem reduces a performance question to the actual pivot
parameter $F$.  Can knot-diagram structure, a braid decomposition, or an
embedded separator order imply a useful bound on $F$ for a specified pivot
policy?  Input sparsity does not suffice.  Nor does good behavior on the
sampled $T(3,q)$ diagrams prove that a degree heuristic always maintains
small neighborhoods.

A first theoretical target is a rigorously specified diagram family for
which the pivot evolution can be described inductively.  A second is a
parameterized fill bound in terms of a checked graph decomposition of the
minor's incidence pattern.  Engineering comparisons should include row-only
sparse elimination, indexed complete pivoting, and dense fallback at the
same field and evaluation points.  Track attempted Schur positions and
transient storage, not only newly created fill.

\subsection{Q7. Can preprocessing expose hidden decomposition?}

Interlacement components characterize visible summands of a fixed diagram.
They do not solve topological prime decomposition.  Can a bounded collection
of Reidemeister or flype-like transformations expose enough decompositions
to bound the largest residual piece on a useful class of inputs?  What
certificate shows that the transformation sequence preserves the knot and
that its search cost is controlled?

One intermediate result would bound the work needed to expose a specified
class of obscured connected sums.  A stronger result would control the
size or width of every remaining factor after a proven reduction strategy.
The algorithm must charge unsuccessful move searches and revalidation,
not just successful crossing reductions.  Memoizing repeated factor PD
codes is useful, but it cannot identify arbitrarily disguised equivalent
summands without further mathematics.

\subsection{Q8. Can a geometric hierarchy meet the state-volume contract?}

Section~\ref{sec:hierarchy} gives a quantitative target:
$\sum_i\log(g_i+1)=O(\log^2 n)$, together with polynomially bounded
epochs, encoded data, and total work per charged visit.  Can an accelerated
hierarchy construction prove that complete contract?  This asks more than
showing termination or a short final certificate.

The first deliverable should be a specification of the state, each
transition, its preconditions, and a verification procedure for its output.
Uniform caps must survive every suffix rebuild.  If a decomposition is
restarted with different caps or a different depth, identify and charge a
new epoch.  Unequal caps and unequal level costs should be retained because
the sharp weighted formula can reveal where work is actually concentrated.
A bound of $O(\log^3 n)$ on the logarithmic state volume would still yield
quasi-polynomial time, but would not be the sharper
$n^{O(\log n)}$ claim.

\subsection{Q9. Can compressed geometry stay compressed through every step?}

Interval orbit counting and binary normal coordinates are useful primitives,
but a chain of routines can lose their advantage at an interface.  Can the
algorithm cut, trace components, transport discs, and rebuild boundary
patterns while retaining compressed multiplicities and explicit bit bounds?
Which routines require an actual list of components, and can that output
be represented by types with multiplicities instead?

Implement a small end-to-end path from a triangulated exterior through one
certified cut and back to a valid handle or triangulation description.
Record input and output bit lengths, represented surface weight, number of
interval identifications, and total primitive calls.  Reject an interface
that silently expands a binary weight into one object per normal disk.
Only after this path is specified should a global hierarchy runtime be
inferred from its building blocks.

\subsection{Q10. What benchmark corpus tests the unresolved decision stage?}

The next corpus should contain growing families of known unknots and
nontrivial knots that survive the complete inexpensive front end.  It should
include alternative diagrams of the same knot, certified move traces for
scrambled unknots, high-width examples, and cases with no visible factors.
Crossing number, effective width, boundary topology, matching support,
object multiplicity, and the deciding stage should all be recorded.

Generated instances need provenance.  A known unknot produced by a verified
move sequence supplies a trusted label; a knot that a cheap invariant
rejects measures a different part of the pipeline.  Separate instances used
for heuristic tuning from those used for evaluation.  Fixed seeds,
randomized paired timing, A/A controls, and retained failures should remain
standard.  The purpose is to identify the remaining mechanism of growth,
not to maximize an average speedup on easily rejected diagrams.

\subsection{Q11. Which certificates should be formalized first?}

The shortest independently verifiable components offer a practical entry
point for proof-assistant work: the interlacement spanning forest and nesting
verifier, the inherited-PD reconstruction, finite-field Schur elimination
with permutation signs, and the mixed-radix ranking and suffix budget.
These have explicit finite data and compact mathematical invariants.

For homology, a later stage can formalize the marked-dot ideal and the
surface-evaluation rules, then Gaussian cancellation and its interaction
with gluing.  The distinction between an isolated-object replay record and
a compact chain-equivalence certificate should remain explicit.  A small
formal checker is more useful than a formal statement whose hardest
topological or complexity premise is hidden in an unimplemented oracle.

\subsection{Proposed order of work}

The first implementation milestone is to integrate the audited factorizer
and sparse filter, preserve the pointed option, and establish the expanded
corpus and counter instrumentation.  The first structural milestone is a
certified interface representation with a proved matching-support bound.
The next algebraic milestone is an exact compositional compression or an
earlier decision certificate that handles a substantial residual class.
In parallel, the hierarchy route needs a small compressed geometric backend
whose complete costs can be charged to the explicit ranking model.

A claim of general quasi-polynomial recognition becomes justified only when
every input is covered by one sound, complete algorithm with a proved
worst-case budget.  Combining fast filters with a conditional scanner and
an unspecified hierarchy subroutine does not yet supply that coverage.
The present package gives concrete improvements and testable mathematical
obligations for the next iteration.
